% Bagian: first-order-logic
% Bab: beyond
% Subbagian: second-order-logic

\documentclass[../../../include/open-logic-section]{subfiles}

\begin{document}

\olfileid[id]{fol}{byd}{sol}

\olsection{Logika Orde Kedua}

Bahasa logika orde kedua memungkinkan kita menguantifikasi bukan hanya atas
!!a{domain} individu, melainkan juga atas relasi pada !!{domain} tersebut.
Diberikan suatu bahasa orde pertama~$\Lang{L}$, untuk setiap $k$ ditambahkan
!!{variable}s~$R$ yang merentang atas relasi beraritas~$k$, dan kuantifikasi atas
variabel-variabel itu diperbolehkan. Jika $R$ merupakan !!a{variable} untuk
suatu relasi beraritas~$k$, dan $t_1$, \dots,~$t_k$ merupakan suku biasa (orde
pertama), maka $\Atom{R}{t_1,\dots,t_k}$ merupakan !!{formula} atomik.
Selain itu, himpunan !!{formula}s didefinisikan tepat seperti dalam logika
orde pertama, dengan klausa tambahan untuk kuantifikasi orde kedua. Perhatikan
bahwa kita hanya mempunyai !!{identity} bagi suku orde pertama: jika $R$ dan
$S$ merupakan !!{variable}s relasi dengan aritas yang sama~$k$, kita dapat
mendefinisikan $\eq[R][S]$ sebagai singkatan bagi
\[
\lforall[x_1 \dots][\lforall[x_k][(\Atom{R}{x_1, \dots, x_k} \liff
  \Atom{S}{x_1, \dots, x_k})]].
\]

Aturan-aturan logika orde kedua sekadar memperluas aturan kuantor ke
variabel-variabel orde kedua yang baru. Akan tetapi, di sini kita harus cukup
cermat dalam menjelaskan bagaimana variabel-variabel ini berinteraksi dengan
!!{predicate}s dari~$\Lang{L}$, serta dengan !!{formula}s dari~$\Lang{L}$
secara lebih umum. Sekurang-kurangnya, variabel relasi berlaku sebagai suku,
sehingga kita mempunyai inferensi berbentuk
\[
!A(R) \Proves \lexists[R][!A(R)]
\]
Namun, jika $\Lang{L}$ adalah bahasa aritmetika dengan simbol relasi
konstan~$<$, kita juga mengharapkan inferensi berikut valid:
\[
x < y \Proves \lexists[R][\Atom{R}{x,y}]
\]
atau, untuk suatu !!{formula}~$!A$,
\[
!A(x_1, \dots, x_k) \Proves \lexists[R][\Atom{R}{x_1,\dots,x_k}]
\]
Secara lebih umum, kita mungkin ingin memperbolehkan inferensi berbentuk
\[
\Subst{!A}{\lambd[\vec x][!B(\vec x)]}{R} \Proves
\lexists[R][!A]
\]
dengan $\Subst{!A}{\lambd[\vec x][!B(\vec x)]}{R}$ menyatakan hasil
penggantian setiap !!{formula} atomik berbentuk
$\Atom{R}{t_1,\dots,t_k}$ dalam~$!A$ dengan $!B(t_1, \dots, t_k)$. Aturan
terakhir ini ekuivalen dengan memiliki suatu {\em skema komprehensi}, yakni
aksioma berbentuk
\[
\lexists[R][\lforall[x_1, \dots, x_k][(!A(x_1, \dots, x_k) \liff
\Atom{R}{x_1, \dots, x_k})]],
\]
satu untuk setiap !!{formula}~$!A$ dalam bahasa orde kedua, dengan $R$ bukan
variabel bebas. (Latihan: tunjukkan bahwa jika $R$ diperbolehkan muncul
dalam~$!A$, skema ini tidak konsisten!)

Ketika ahli logika merujuk pada ``aksioma logika orde kedua'', yang biasanya
dimaksud ialah perluasan minimal logika orde pertama dengan aturan kuantor
orde kedua beserta skema komprehensi. Namun, sering kali menarik untuk
mengkaji subsistem yang lebih lemah dari aksioma dan aturan tersebut. Sebagai
contoh, perhatikan bahwa dalam bentuknya yang sepenuhnya umum, skema aksioma
komprehensi bersifat \emph{impredikatif}: skema itu memungkinkan kita
menyatakan keberadaan suatu relasi $\Atom{R}{x_1, \dots, x_k}$ yang
``didefinisikan'' oleh !!a{formula} dengan kuantor orde kedua; dan kuantor ini
merentang atas himpunan semua relasi semacam itu---himpunan yang mencakup $R$
sendiri!{} Menjelang pergantian abad kedua puluh, salah satu tanggapan yang
umum terhadap paradoks Russell ialah menyalahkan definisi semacam itu dan
menghindarinya dalam pengembangan landasan matematika. Jika penggunaan
kuantor orde kedua dalam !!{formula}~$!A$ dilarang, diperoleh bentuk
komprehensi {\em predikatif}, yang agak lebih lemah.

Dari sudut pandang semantis, suatu !!{structure} orde kedua dapat dipandang
sebagai terdiri atas suatu !!{structure} orde pertama untuk bahasanya, yang
dipasangkan dengan suatu himpunan relasi pada !!{domain} yang menjadi rentang
kuantor orde kedua (lebih tepatnya, untuk setiap $k$ terdapat suatu himpunan
relasi beraritas~$k$). Tentu saja, jika komprehensi disertakan dalam sistem
!!{derivation}, terdapat persyaratan tambahan bahwa tersedia cukup banyak
relasi dalam ``bagian orde kedua'' untuk memenuhi aksioma komprehensi---jika
tidak, sistem !!{derivation} itu tidak sahih!{} Salah satu cara mudah untuk
memastikan bahwa tersedia cukup banyak relasi ialah mengambil bagian orde
kedua sebagai terdiri atas \emph{semua} relasi pada bagian orde pertama.
Kita menyebut !!a{structure} semacam itu \emph{penuh}, dan dalam arti tertentu
benar-benar merupakan ``!!{structure} yang dimaksudkan'' bagi bahasa tersebut.
Jika perhatian kita dibatasi pada !!{structure}s penuh, kita memperoleh apa
yang dikenal sebagai semantik orde kedua \emph{penuh}. Dalam hal itu,
menentukan suatu !!{structure} bermuara pada menentukan bagian orde pertamanya,
karena isi bagian orde kedua mengikuti secara implisit dari bagian itu.

Ringkasnya, terdapat sejumlah ambiguitas ketika membicarakan logika orde
kedua. Dari segi sistem !!{derivation}, yang dimaksud dapat berupa salah satu
dari
\begin{enumerate}
\item Sistem !!{derivation} orde kedua ``minimal'', beserta beberapa aksioma
  komprehensi.
\item Sistem !!{derivation} orde kedua ``standar'', dengan komprehensi penuh.
\end{enumerate}
Dari segi semantik, yang diminati dapat berupa salah satu dari
\begin{enumerate}
\item Semantik ``lemah'', dengan !!a{structure} terdiri atas bagian orde
  pertama beserta bagian orde kedua yang cukup besar untuk memenuhi aksioma
  komprehensi.
\item Semantik orde kedua ``standar'', yang hanya mempertimbangkan
  !!{structure}s penuh.
\end{enumerate}
Ketika ahli logika tidak menentukan sistem !!{derivation} atau semantik yang
dimaksud, biasanya mereka merujuk pada butir kedua dalam masing-masing daftar.
Keuntungan menggunakan semantik ini ialah, sebagaimana akan kita lihat, bahwa
semantik ini memberikan deskripsi kategoris bagi banyak struktur matematika
alami; sementara itu, sistem !!{derivation}-nya cukup kuat dan sahih bagi
semantik ini. Kekurangannya ialah bahwa sistem !!{derivation} tersebut
\emph{tidak} lengkap bagi semantik itu; bahkan, \emph{tidak ada} sistem
!!{derivation} yang diberikan secara efektif yang lengkap bagi semantik orde
kedua penuh. Di sisi lain, kita akan melihat bahwa sistem !!{derivation}
tersebut \emph{memang} lengkap bagi semantik yang diperlemah; ini mengakibatkan bahwa
jika suatu kalimat tidak dapat dibuktikan, maka terdapat \emph{suatu}
!!{structure}, yang tidak harus penuh, tempat kalimat itu salah.

Bahasa logika orde kedua sangat kaya. Relasi uner dapat diidentifikasi dengan
subhimpunan dari !!{domain}, sehingga kita khususnya dapat menguantifikasi
atas himpunan-himpunan ini; misalnya, induksi bagi bilangan asli dapat
diungkapkan dengan satu aksioma
\[
\lforall[R][((\Atom{R}{\Obj{0}} \land \lforall[x][(\Atom{R}{x} \lif
    \Atom{R}{x'})]) \lif \lforall[x][\Atom{R}{x}])].
\]
Jika bahasa aritmetika diambil mempunyai simbol $\Obj 0, \Obj \prime, +,
\times$, dan $<$, aksioma-aksioma berikut dapat ditambahkan untuk menguraikan
perilakunya:
\begin{enumerate}
\item $\lforall[x][\lnot x' = \Obj 0]$
\item $\lforall[x][\lforall[y][(x' = y' \lif x = y)]]$
\item $\lforall[x][(x + \Obj 0) = x]$
\item $\lforall[x][\lforall[y][(x + y') = (x + y)']]$
\item $\lforall[x][(x \times \Obj 0) = \Obj 0]$
\item $\lforall[x][\lforall[y][(x \times y') = ((x \times y) + x)]]$
\item $\lforall[x][\lforall[y][(x < y \liff \lexists[z][y = (x + z')])]]$
\end{enumerate}
Tidak sulit untuk menunjukkan bahwa aksioma-aksioma ini, bersama aksioma
induksi di atas, memberikan deskripsi kategoris bagi !!{structure}~$\Struct{N}$,
model standar aritmetika, asalkan kita menggunakan semantik orde kedua penuh.
Diberikan sembarang !!{structure}~$\Struct{M}$ tempat aksioma-aksioma ini
benar, definisikan suatu fungsi~$f$ dari $\Nat$ ke !!{domain} dari $\Struct{M}$
dengan rekursi biasa pada $\Nat$, sedemikian sehingga $f(0) =
\Assign{\Obj 0}{M}$ dan $f(x+1) = \Assign{\prime}{M}(f(x))$. Dengan
menggunakan induksi biasa pada~$\Nat$ dan fakta bahwa aksioma (1) dan~(2)
berlaku dalam~$\Struct M$, kita melihat bahwa $f$ bersifat !!{injective}.
Untuk melihat bahwa $f$ bersifat !!{surjective}, misalkan~$P$ adalah himpunan
elemen-elemen dari~$\Domain{M}$ yang berada dalam daerah hasil~$f$. Karena
$\Struct M$ penuh, $P$ berada dalam !!{domain} orde kedua. Berdasarkan
konstruksi~$f$, kita mengetahui bahwa $\Assign{\Obj 0}{M}$ berada dalam~$P$,
dan bahwa $P$~tertutup terhadap $\Assign{\prime}{M}$. Fakta bahwa aksioma
induksi berlaku dalam $\Struct M$ (khususnya, bagi~$P$) menjamin bahwa
$P$~sama dengan seluruh !!{domain} orde pertama dari~$\Struct M$. Ini
menunjukkan bahwa $f$ merupakan !!a{bijection}. Menunjukkan bahwa $f$ adalah
homomorfisme juga tidak lebih sulit, dengan menggunakan induksi biasa pada
$\Nat$ berulang kali.

Dalam istilah teori himpunan, suatu fungsi hanyalah jenis relasi khusus;
misalnya, fungsi uner~$f$ dapat diidentifikasi dengan relasi biner~$R$ yang
memenuhi $\lforall[x][\lexists![y][R(x,y)]]$. Akibatnya, kita juga dapat
menguantifikasi atas fungsi. Dengan menggunakan semantik penuh, kita kemudian
dapat mendefinisikan kelas !!{structure}s tak hingga sebagai kelas semua
!!{structure}s~$\Struct M$ yang memiliki fungsi injektif dari !!{domain}
$\Struct M$ ke suatu subhimpunan sejati dari domain itu sendiri:
\[
\lexists[f][(\lforall[x][\lforall[y][(\eq[f(x)][f(y)] \lif \eq[x][y])]]
  \land \lexists[y][\lforall[x][\eq/[f(x)][y]]])].
\]
Negasi kalimat ini kemudian mendefinisikan kelas !!{structure}s berhingga.

Selain itu, kita dapat mendefinisikan kelas urutan-baik dengan menambahkan
hal berikut pada definisi urutan linear:
\[
\lforall[P][(\lexists[x][\Atom{P}{x}] \lif \lexists[x][(\Atom{P}{x}
    \land \lforall[y][(y < x \lif \lnot \Atom{P}{y})])])].
\]
Ini menyatakan bahwa setiap himpunan tak kosong mempunyai elemen terkecil,
dengan mengidentifikasi ``himpunan'' dengan ``relasi satu-tempat''. Sebagai
contoh lain, gagasan keterhubungan bagi graf dapat diungkapkan dengan
menyatakan bahwa tidak ada pemisahan nontrivial atas verteks menjadi
bagian-bagian yang tidak terhubung:
\[
\lnot \lexists[A][(\lexists[x][A(x)] \land \lexists[y][\lnot A(y)]
  \land \lforall[w][\lforall[z][((\Atom{A}{w} \land \lnot \Atom{A}{z})
      \lif \lnot \Atom{R}{w,z})]])].
\]
Sebagai contoh lain lagi, cobalah sebagai latihan mendefinisikan kelas
!!{structure}s berhingga yang !!{domain}-nya berukuran genap. Yang lebih
mengagumkan, kita dapat memberikan deskripsi kategoris bagi bilangan real
sebagai lapangan terurut lengkap yang memuat bilangan rasional.

Singkatnya, logika orde kedua jauh lebih ekspresif daripada logika orde
pertama. Itulah kabar baiknya; sekarang kabar buruknya. Kita telah menyebutkan
bahwa tidak ada sistem !!{derivation} efektif yang lengkap bagi semantik orde
kedua penuh. Baik disukai maupun tidak, banyak sifat logika orde pertama tidak
lagi berlaku, termasuk kekompakan dan teorema L\"owenheim--Skolem.

Di sisi lain, jika kita bersedia melepaskan semantik orde kedua penuh dan
menggunakan semantik yang lebih lemah, sistem !!{derivation} orde kedua
minimal lengkap bagi semantik ini. Dengan kata lain, jika $\Proves$ menyatakan
keterbuktian dalam sistem minimal dan $\Entails$ menyatakan perikutan dalam
semantik yang lebih lemah, kita dapat menunjukkan bahwa setiap kali $\Gamma
\Entails !A$, berlaku $\Gamma \Proves !A$. Jika aksioma komprehensi tertentu
hendak disertakan dalam sistem
!!{derivation}, semantik harus dibatasi pada struktur orde kedua yang memenuhi
aksioma-aksioma tersebut: misalnya, jika $\Delta$ terdiri atas suatu himpunan
aksioma komprehensi (mungkin semuanya), kita memperoleh bahwa jika $\Gamma
\cup \Delta \Entails !A$, maka $\Gamma \cup \Delta \Proves !A$. Khususnya,
jika $!A$ tidak dapat dibuktikan dengan menggunakan aksioma komprehensi yang
sedang dipertimbangkan, maka terdapat model dari $\lnot !A$ tempat
aksioma-aksioma komprehensi tersebut tetap berlaku.

Cara termudah untuk melihat bahwa teorema kelengkapan berlaku bagi semantik
yang lebih lemah ialah memandang logika orde kedua sebagai logika banyak-sorta
sebagai berikut. Satu sorta ditafsirkan sebagai !!{domain} ``orde pertama''
biasa, lalu untuk setiap $k$ terdapat !!a{domain} ``relasi beraritas~$k$''.
Bahasanya diambil mempunyai simbol relasi bawaan
``$\Atom{\Obj{true}_k}{R,x_1,\dots,x_k}$'', yang dimaksudkan untuk menyatakan
bahwa $R$ berlaku pada $x_1$, \dots,~$x_k$, dengan $R$ variabel dari sorta
``relasi beraritas~$k$'' dan $x_1$, \dots,~$x_k$ objek dari sorta orde pertama.

Dengan identifikasi ini, semantik orde kedua lemah pada hakikatnya merupakan
semantik biasa bagi logika banyak-sorta; dan kita telah mengamati bahwa logika
banyak-sorta dapat ditanamkan ke dalam logika orde pertama. Jadi, dengan
memperhitungkan penerjemahan bolak-balik tersebut, konsepsi logika orde kedua
yang lebih lemah sebenarnya merupakan suatu bentuk logika orde pertama yang
menyamar, dengan !!{domain} memuat baik ``objek'' maupun ``relasi'' yang
diatur oleh aksioma yang sesuai.

\end{document}
